
\documentclass[11pt,a4paper]{article}

\usepackage[T1]{fontenc}
\usepackage[utf8]{inputenc}
\usepackage[a4paper,margin=1in]{geometry}
\usepackage{lmodern}
\usepackage{microtype}
\usepackage{amsmath,amssymb,mathtools,bm}
\usepackage{booktabs,longtable,array}
\usepackage{xcolor}
\usepackage{enumitem}
\usepackage{listings}
\usepackage{hyperref}

\hypersetup{
  colorlinks=true,
  linkcolor=black,
  citecolor=black,
  urlcolor=black,
  pdftitle={What Matter's Most: Finding Time In A temporal HyperSpherical Zenith},
  pdfauthor={Daniel Alexander Trawin},
  pdfcreator={Trawin, Daniel Alexander},
  pdfproducer={Trawin, Daniel Alexander},
  pdfsubject={Hyperspherical resonance, Bessel half-mode drops, temporal accumulation, gravitational charge, and residual curvature},
  pdfkeywords={TZPID, Trawin, hyperspherical resonance, Bessel equation, entropy fold, temporal accumulation, gravitational charge, matter creation, accumulated force, DOI 10.5281/zenodo.20579476, version timestamp 2026-06-07 07:04:10 -04:00}
}

\pdfinfo{
  /Title (What Matter's Most: Finding Time In A temporal HyperSpherical Zenith)
  /Author (Daniel Alexander Trawin)
  /Creator (Trawin, Daniel Alexander)
  /Subject (Hyperspherical resonance, Bessel half-mode drops, temporal accumulation, gravitational charge, and residual curvature)
  /Keywords (TZPID, Trawin, hyperspherical resonance, Bessel equation, entropy fold, temporal accumulation, gravitational charge, matter creation, accumulated force, DOI 10.5281/zenodo.20579476, version timestamp 2026-06-07 07:04:10 -04:00)
}

\lstdefinestyle{proofcode}{
  basicstyle=\ttfamily\footnotesize\color{black},
  keywordstyle=\color{black},
  commentstyle=\color{black},
  stringstyle=\color{black},
  identifierstyle=\color{black},
  breaklines=true,
  columns=fullflexible,
  keepspaces=true,
  frame=none,
  rulecolor=\color{black},
  showstringspaces=false
}
\lstset{style=proofcode}

\renewcommand{\boxed}[1]{#1}
\renewcommand{\fbox}[1]{#1}

\newcommand{\dd}{\,\mathrm{d}}
\newcommand{\R}{\mathbb{R}}
\newcommand{\Sph}{\mathbb{S}}
\newcommand{\half}{\frac{1}{2}}
\newcommand{\Teff}{T_{\mu\nu}^{\mathrm{eff}}}
\newcommand{\Tmat}{T_{\mu\nu}^{\mathrm{matter}}}
\newcommand{\Tsound}{T_{\mu\nu}^{\mathrm{sound}}}
\newcommand{\Sigmahalf}{\Sigma_{\mu\nu}^{(1/2)}}

\title{\textbf{What Matter's Most: Finding Time In A temporal HyperSpherical Zenith}}
\author{Daniel Alexander Trawin\\
\small DOI: \href{https://doi.org/10.5281/zenodo.20579476}{10.5281/zenodo.20579476}}
\date{\today\\\small Version timestamp: 2026-06-07 07:04:10 -04:00}

\begin{document}
\maketitle

\begin{abstract}
This paper develops a sharpened version of the thesis that ``sound'' can be made into the engine of a geometric matter-and-gravity model, provided that sound is not interpreted as ordinary acoustic pressure in air. Instead, sound is treated as a quantized resonant compression or phase mode: a phonon-like vibration of a field, condensate, plasma, or hyperspherical substrate. The proof spine combines four mutually reinforcing signatures: hyperspherical enclosure acoustics, Bessel boundary quantization, cross-scale ripple ratios, and Pythagorean/Hopf reciprocal holonomy. The resulting chain is
\begin{center}
\fbox{\begin{minipage}{0.88\linewidth}
\centering
hyperspherical resonance $\rightarrow$ disharmonic Bessel fold $\rightarrow$ entropy-producing half-mode drop $\rightarrow$ temporal accumulation $\rightarrow$ curvature, time-arrow, and gravitational charge
\end{minipage}}
\end{center}
The resulting model is not divided into a merely solid accounting layer and a separate conjectural layer. It is a unified validation ladder: standard mass-energy accounting fixes the Planck-scaled source; enclosure acoustics fixes the allowed modes; Bessel zeros fix the boundary spectrum; ripple ratios and Pythagorean holonomy supply cross-domain projection tests; and the residual target asks whether the accumulated source contains a measurable half-Bessel signature beyond ordinary mass-energy.
\end{abstract}

\begingroup
\setcounter{tocdepth}{1}
\scriptsize
\renewcommand{\baselinestretch}{0.76}\selectfont
\setlength{\parskip}{0pt}
\tableofcontents
\endgroup
\newpage

\section{Sound as a Field-Resonance Engine}

The proposal begins with a clarification. The engine of the model is not ordinary sound in air. Ordinary acoustic pressure requires a material medium. The relevant object is sound in the generalized sense: a quantized resonant compression or phase mode of a field, condensate, plasma, or hyperspherical substrate. In this form, ``sound'' can be written as a seed mechanism for curvature and matter formation.

The sharpened thesis is the following sequence:
\begin{equation}
\boxed{
\begin{gathered}
\text{hyperspherical resonance}
\rightarrow
\text{disharmonic Bessel fold}
\rightarrow
\text{entropy-producing half-mode drop}
\\
\rightarrow
\text{temporal accumulation}
\rightarrow
\text{curvature, time-arrow, and gravitational charge}.
\end{gathered}}
\label{eq:thesis-chain}
\end{equation}
The model is therefore not merely an acoustic analogy. It is a proposed causal spine: resonant field modes generate a half-mode defect; the defect produces entropy; entropy fixes a time direction; and a causal accumulation kernel turns the resulting effective stress into curvature.

\section{Hyperspherical Enclosure and Projection Semantics}

The semantic core of the model is a nested hyperspherical enclosure. The observable three-dimensional pattern is treated as a projection of standing modes defined on, or within, a higher-dimensional spherical boundary. In its simplest angular-radial form, the mode basis is
\begin{equation}
\Phi(r,\theta,\phi,t)=A_{\ell m q}j_\ell(k_{\ell q}r)Y_{\ell m}(\theta,\phi)e^{-i\omega t}.
\label{eq:spherical-mode-basis}
\end{equation}
Here $Y_{\ell m}$ carries the angular enclosure mode, while $j_\ell(k r)$ carries the radial standing wave. The allowed $k$ values are not arbitrary; they are selected by boundary conditions at the enclosure radius.

This same mathematical structure appears in the cosmological acoustic record. The baryon acoustic oscillation term can be written schematically as
\begin{equation}
\delta_b(k)=A\,j_0(k r_s)T(k),
\label{eq:bao-bessel}
\end{equation}
where $r_s$ is the sound horizon and $T(k)$ is the transfer function. The zeros of $j_0(k r_s)$ occur at
\begin{equation}
k_n r_s=n\pi,
\label{eq:bao-nodes}
\end{equation}
which is the same root condition that selects the normal modes of a spherical acoustic cavity. The observed matter power spectrum,
\begin{equation}
P(k)=P_{\mathrm{prim}}(k)T^2(k),
\label{eq:power-spectrum}
\end{equation}
therefore contains a fossilized enclosure-acoustic spacing. In this reading, the filament web is not merely analogous to spherical acoustics; its characteristic scale is selected by the same eigenvalue problem.

The projection layer is expressed by maps
\begin{equation}
\pi_i:\mathcal{T}\longrightarrow\mathbb{R}^3,
\label{eq:projection-map}
\end{equation}
where $\mathcal{T}$ denotes the nested hyperspherical source structure and $\pi_i$ denotes a projection into an observable three-dimensional domain. A dimensionless ratio that survives this projection is more important than an absolute size, because projection can rescale length while preserving mode ratios. This is the reason ripple shape ratios, musical frequency ratios, and cosmic acoustic scales can belong to one proof spine without being the same physical medium.

\section{The Bessel Starting Point}

The ordinary Bessel equation is
\begin{equation}
 x^2\frac{\dd^2 y}{\dd x^2}+x\frac{\dd y}{\dd x}+(x^2-n^2)y=0.
\label{eq:bessel-n}
\end{equation}
Replacing the integer order $n$ by a generalized hyperspherical order $\nu$ gives
\begin{equation}
\boxed{x^2y''+xy'+(x^2-\nu^2)y=0.}
\label{eq:bessel-nu}
\end{equation}
The standard solution basis is
\begin{equation}
 y(x)=A J_\nu(x)+B Y_\nu(x),
\label{eq:bessel-solutions}
\end{equation}
where $J_\nu$ and $Y_\nu$ are Bessel functions of the first and second kind.

Now place this mode inside a hyperspherical enclosure of radius $R$. Let the resonant field separate as
\begin{equation}
 \Phi(r,\Omega,t)=R_\ell(r)Y_{\ell\mathbf m}(\Omega)e^{-i\omega t}.
\label{eq:field-separation}
\end{equation}
For a $d$-dimensional hyperspherical enclosure, the radial wave equation is
\begin{equation}
 r^2R_\ell''+(d-1)rR_\ell'+\left(k^2r^2-\ell(\ell+d-2)\right)R_\ell=0.
\label{eq:d-radial}
\end{equation}
Make the substitution
\begin{equation}
 x=kr,\qquad R_\ell(r)=r^{-(d-2)/2}y(x).
\label{eq:radial-substitution}
\end{equation}
The radial equation becomes
\begin{equation}
\boxed{x^2y''+xy'+\left(x^2-\nu_d^2\right)y=0,}
\label{eq:hyperspherical-bessel}
\end{equation}
where
\begin{equation}
\boxed{\nu_d=\ell+\frac{d-2}{2}.}
\label{eq:nu-d}
\end{equation}
For a true four-spatial-dimensional hyperspherical enclosure, namely a four-ball with a three-sphere boundary, $d=4$ and therefore
\begin{equation}
\boxed{\nu_4=\ell+1.}
\label{eq:nu-four}
\end{equation}
Thus the hyperspherical Bessel radial mode may be written as
\begin{equation}
\boxed{R_\ell(r)=\frac{1}{r}J_{\ell+1}(kr).}
\label{eq:radial-four}
\end{equation}
Enclosure quantization is imposed by the boundary condition
\begin{equation}
\boxed{J_{\ell+1}(kR)=0.}
\label{eq:boundary-zero}
\end{equation}
Equivalently,
\begin{equation}
\boxed{k_{\ell q}=\frac{j_{\ell+1,q}}{R},}
\label{eq:k-zeros}
\end{equation}
where $j_{\ell+1,q}$ is the $q$-th zero of $J_{\ell+1}$.

This is the first clean translation of the visual intuition into a hyperspherical sound-resonance enclosure: allowed modes are not arbitrary; they are selected by Bessel zeros.

\section{Disharmonic Fold and the Half-Bessel Drop}

A harmonic enclosure holds clean standing modes,
\begin{equation}
 \omega_{\ell q}=v_s k_{\ell q}=v_s\frac{j_{\nu,q}}{R}.
\label{eq:harmonic-mode}
\end{equation}
Here $v_s$ is the wave speed. For ordinary sound, it is the speed of sound in the medium; for a vacuum-field version, it is replaced by an effective propagation speed, often $c$ or a model-specific phase velocity.

A disharmonic fold is defined as a phase-inverted half-order deformation of the Bessel mode:
\begin{equation}
\boxed{\mathfrak{F}_{1/2}:J_\nu(x)\longrightarrow J_{\nu-1/2}(x+\delta x)e^{i\pi}.}
\label{eq:fold-map}
\end{equation}
Because $e^{i\pi}=-1$, the folded mode is phase-opposed, not merely shifted. The fold field may be written
\begin{equation}
\boxed{
\Phi_{\mathrm{fold}}
=
A J_\nu(kr)e^{-i\omega t}
-
\eta A J_{\nu-1/2}((k+\delta k)r)e^{-i(\omega+\delta\omega)t}.
}
\label{eq:fold-field}
\end{equation}
This is the disharmonic fold: the original resonance and the half-dropped resonance interfere destructively and constructively across the hyperspherical cavity.

The allowed modes are fixed by
\begin{equation}
 k_{\nu q}=\frac{j_{\nu,q}}{R}.
\label{eq:k-nu-q}
\end{equation}
A half-Bessel drop is
\begin{equation}
\boxed{\nu\rightarrow\nu-\frac12.}
\label{eq:half-drop}
\end{equation}
Thus the resonant wavenumber drops from $k_{\nu q}$ to
\begin{equation}
 k_{\nu-1/2,q}=\frac{j_{\nu-1/2,q}}{R}.
\label{eq:k-half}
\end{equation}
The released or reconfigured mode energy is
\begin{equation}
\boxed{
\Delta E^{(1/2)}_{\nu q}=\hbar v_s\left(\frac{j_{\nu,q}-j_{\nu-1/2,q}}{R}\right).
}
\label{eq:delta-e-half}
\end{equation}

For the fundamental hyperspherical case,
\begin{equation}
 d=4,\qquad \ell=0,\qquad \nu=1.
\label{eq:fundamental-case}
\end{equation}
The first zero of $J_1$ is
\begin{equation}
 j_{1,1}\approx 3.83170597.
\label{eq:j11}
\end{equation}
Meanwhile,
\begin{equation}
 J_{1/2}(x)\propto\frac{\sin x}{\sqrt{x}},
\label{eq:j-half-form}
\end{equation}
so its first zero is
\begin{equation}
 j_{1/2,1}=\pi\approx 3.14159265.
\label{eq:j-half-zero}
\end{equation}
The fractional half-Bessel drop is therefore
\begin{equation}
\boxed{
\delta_{B,1/2}=\frac{j_{1,1}-j_{1/2,1}}{j_{1,1}}\approx 0.1801060212.
}
\label{eq:half-drop-fingerprint}
\end{equation}
The lowest hyperspherical resonance therefore produces an approximately $18.01\%$ mode drop. This is not by itself proof of the physical theory. It is, however, a concrete mathematical fingerprint. The fractional mode drop $\delta_{B,1/2}\approx 0.180106$ has been computationally checked in the Wolfram Language; the corresponding proof certificate is recorded in Appendix~\ref{app:formal-verification}.

\section{Cross-Scale Ripple Ratios}

The ripple evidence supplies the spatial version of the projection test. A ripple is meaningful here not because a sand ripple, a water ripple, and a sky or plasma ripple share the same absolute wavelength, but because they preserve dimensionless shape information across changes of size and medium. The relevant empirical quantity is the ripple index
\begin{equation}
\mathrm{RI}=\frac{\lambda}{h},
\label{eq:ripple-index}
\end{equation}
where $\lambda$ is wavelength and $h$ is ripple height.

Measured ripple indices form a regime-dependent but scale-free family:
\begin{equation}
\mathrm{RI}_{\mathrm{sand}}\approx 15\text{--}20,\qquad
\mathrm{RI}_{\mathrm{water}}\approx 4\text{--}16,\qquad
\mathrm{RI}_{\mathrm{current}}\approx 8\text{--}20.
\label{eq:ripple-index-family}
\end{equation}
The point is not that every medium has one universal spatial number. The point is that the observable pattern is organized by dimensionless ratios. Sand ripples, megaripples, and dunes can differ by orders of magnitude in length while remaining part of a self-similar bedform hierarchy. In the enclosure picture, this is what projection should do: preserve ratios while allowing the projected length scale to vary.

This also clarifies the role of the registry ratio $32/27\approx 1.18518$. It is not a spatial ripple index. It is the Pythagorean minor third, a frequency ratio. Thus the ripple spine has two axes:
\begin{align}
\text{spatial axis:}\quad & \mathrm{RI}=\lambda/h,\\
\text{frequency axis:}\quad & r_{\mathrm{minor\ third}}=\frac{32}{27}.
\end{align}
The spatial axis is already supported by measured ripple-index families. The frequency axis is a testable mode-selection prediction for cymatic, Faraday, or spherical-enclosure pattern experiments.

\section{Pythagorean Interval, Comma, and Reciprocal Holonomy}

The Pythagorean bridge should be stated with the music theory exact. The ratio $3/2$ is the perfect fifth, not the third. The Pythagorean major third is $81/64$, and the Pythagorean minor third is $32/27$. The circle-of-fifths mechanism, however, is generated by repeated multiplication by $3/2$:
\begin{equation}
\gamma
=
\frac{(3/2)^{12}}{2^7}
=
\frac{3^{12}}{2^{19}}
=
\frac{531441}{524288}
\approx 1.0136433.
\label{eq:pythagorean-comma}
\end{equation}
This is the Pythagorean comma. In cents,
\begin{equation}
c_\gamma=1200\log_2(\gamma)\approx 23.460.
\label{eq:comma-cents}
\end{equation}
The circle of fifths returns to the same pitch class, but not to the same pitch. The residual is a non-closing loop, and a residual produced by carrying a state around a closed loop is a holonomy.

On an octave-reduced pitch circle, the residual rotation is
\begin{equation}
\theta_\gamma
=
2\pi\left(12\log_2\frac32-7\right)
\approx 0.122836\ \mathrm{rad}.
\label{eq:comma-angle}
\end{equation}
In the hyperspherical enclosure reading, this is the audible boundary trace of a bulk projection. A loop on the base $S^2$ of a Hopf fibration lifts into the bulk $S^3$ with a fiber rotation
\begin{equation}
\Delta\chi_{\mathrm{Hopf}}=\Omega=\theta_\gamma,
\label{eq:hopf-phase}
\end{equation}
and the reciprocal outward flip is
\begin{equation}
\omega_{\mathrm{bulk}}=\gamma^{-1}=\frac{524288}{531441}.
\label{eq:bulk-reciprocal}
\end{equation}
Thus the same map appears in multiple domains:
\begin{equation}
r\longmapsto r^{-1},
\label{eq:reciprocal-map}
\end{equation}
as fifth to descending fifth, comma excess to bulk deficit, and avalanche exponent $3/2$ to cascade reciprocal $2/3$. The Pythagorean part of the proof spine is therefore not ornamental. It is the frequency-ratio expression of the same projection logic that the ripple index expresses spatially.

\section{Entropy Correction and the Direction of Time}

Let the half-Bessel drop create an entropy change
\begin{equation}
 \Delta S_{1/2}=k_B\ln\left(\frac{\Omega_{\nu-1/2,q}}{\Omega_{\nu,q}}\right),
\label{eq:delta-s}
\end{equation}
where $\Omega$ counts accessible microstates of the mode.

The entropy-gradient curvature seed is
\begin{equation}
\boxed{
\Sigmahalf
=
\ell_P^2
\left(\nabla_\mu\nabla_\nu-g_{\mu\nu}\Box\right)
\left(\frac{\Delta S_{1/2}}{k_B}\right).
}
\label{eq:sigma-half}
\end{equation}
The Planck length is
\begin{equation}
 \ell_P=\sqrt{\frac{\hbar G}{c^3}}.
\label{eq:planck-length}
\end{equation}
Define an effective source tensor by
\begin{equation}
\boxed{
\Teff
=
\Tmat+
\Tsound+
\frac{c^4}{8\pi G\ell_P^2}\Sigmahalf.
}
\label{eq:t-eff}
\end{equation}
In plain terms,
\begin{equation}
\boxed{
\text{matter stress}
+
\text{sound/resonance stress}
+
\text{entropy-fold stress}
=
\text{effective gravitational source}.
}
\label{eq:source-plain}
\end{equation}
The half-Bessel drop does not merely move energy around. It creates an irreversible entropy gradient. In this model, that entropy gradient gives the arrow of time.

The causation chain is
\begin{equation}
\boxed{
\Delta E_{\mathrm{Bessel}}
\rightarrow \Delta S\rightarrow
\text{temporal ordering}\rightarrow\text{curvature accumulation}.
}
\label{eq:entropy-causation}
\end{equation}

\section{Temporal Accumulation Kernel}

Use a normalized temporal accumulation kernel
\begin{equation}
\boxed{
\mathcal{K}(t,\tau)=\frac{1}{\tau_{\mathrm{dec}}}e^{-(t-\tau)/\tau_{\mathrm{dec}}}\Theta(t-\tau).
}
\label{eq:kernel}
\end{equation}
The factor $\Theta(t-\tau)$ enforces causal ordering: only the past contributes to the present. The normalization keeps the dimensions controlled.

The accumulated curvature tensor is then
\begin{equation}
\boxed{
\mathcal{G}_{\mu\nu}^{\mathrm{acc}}(x,t)
=
\frac{8\pi G}{c^4}
\int_{-\infty}^{t}\mathcal{K}(t,\tau)\Teff(x,\tau)\dd\tau.
}
\label{eq:g-acc}
\end{equation}
Substituting the sound-fold-entropy source gives
\begin{equation}
\boxed{
\mathcal{G}_{\mu\nu}^{\mathrm{acc}}
=
\frac{8\pi G}{c^4}
\int_{-\infty}^{t}\mathcal{K}(t,\tau)
\left[
\Tmat+
\Tsound+
\frac{c^4}{8\pi G\ell_P^2}\Sigmahalf
\right]\dd\tau.
}
\label{eq:g-acc-full}
\end{equation}
Equivalently,
\begin{align}
\boxed{
\mathcal{G}_{\mu\nu}^{\mathrm{acc}}
}&=\frac{8\pi G}{c^4}\int_{-\infty}^{t}\mathcal{K}\Tmat\dd\tau
+\frac{8\pi G}{c^4}\int_{-\infty}^{t}\mathcal{K}\Tsound\dd\tau
\notag\\
&\quad+\frac{1}{\ell_P^2}\int_{-\infty}^{t}\mathcal{K}\Sigmahalf\dd\tau.
\label{eq:g-acc-decomp}
\end{align}
The concise interpretation is
\begin{center}
\fbox{\begin{minipage}{0.88\linewidth}
\centering
\textbf{curvature}
= accumulated matter
+ accumulated sound
+ Planck-scaled accumulated entropy-fold correction.
\end{minipage}}
\end{center}
This is the cleanest version of the idea: curvature is accumulated matter plus accumulated resonance plus a Planck-scaled entropy-fold residual.
The kernel normalization and residual decomposition in this section are also part of the formal verification layer: Wolfram checks the analytic reductions, while Isabelle records the bounded logical chain from large-number smoothing to effective stress-energy and accumulated curvature. The scripts are collected in Appendix~\ref{app:formal-verification}.

\section{Temporal Frame-Dragging}

Frame-dragging is associated with the off-diagonal time-space pieces of the metric, often denoted $g_{0i}$ or $h_{0i}$ in weak-field form. In the accumulated model, the source is the accumulated energy-current
\begin{equation}
 J_i^{\mathrm{eff}}=\frac{T_{0i}^{\mathrm{eff}}}{c}.
\label{eq:j-eff}
\end{equation}
Inside a hyperspherical enclosure, write
\begin{equation}
\boxed{
 h_{0a}^{\mathrm{acc}}(X,t)
 =
 -\frac{4G}{c^3}
 \int_{\Sph_R^3}\dd V'\,G_R(X,X')
 \int_{-\infty}^{t}\mathcal{K}(t,\tau)J_a^{\mathrm{eff}}(X',\tau)\dd\tau.
}
\label{eq:frame-dragging}
\end{equation}
Here $X\in\Sph_R^3$ is a point on or within the hyperspherical enclosure, $G_R(X,X')$ is the hyperspherical Green function, and $a$ is a hyperspherical angular direction.

Thus the phrase ``time frame-dragging'' receives the equation-level reading
\begin{equation}
\boxed{
\text{time frame-dragging}
=
\text{temporally accumulated circulating sound/mass/entropy current}.
}
\label{eq:time-frame-dragging-plain}
\end{equation}

\section{Planck-Scaled Gravitational Charge}

The exact gravitational charge of a particle can be made exact only after defining what ``gravitational charge'' means. The cleanest Planck-scaled definition is
\begin{equation}
\boxed{q_g(X)=\frac{m_X}{m_P},}
\label{eq:qg}
\end{equation}
where
\begin{equation}
\boxed{m_P=\sqrt{\frac{\hbar c}{G}}.}
\label{eq:planck-mass}
\end{equation}
Equivalently,
\begin{equation}
\boxed{q_g(X)=m_X\sqrt{\frac{G}{\hbar c}}.}
\label{eq:qg-equivalent}
\end{equation}
The gravitational coupling between two particles becomes
\begin{equation}
\boxed{
\alpha_G(X,Y)=\frac{Gm_Xm_Y}{\hbar c}=q_g(X)q_g(Y).
}
\label{eq:alpha-g}
\end{equation}
This definition makes gravitational charge dimensionless.

Using the Planck-scaled charge, the indicative single-particle values are:
\begin{equation}
\boxed{q_g(p)\approx 7.68515\times 10^{-20},}
\label{eq:q-proton}
\end{equation}
\begin{equation}
\boxed{q_g(n)\approx 7.69574\times 10^{-20},}
\label{eq:q-neutron}
\end{equation}
\begin{equation}
\boxed{q_g(e)\approx 4.18546\times 10^{-23}.}
\label{eq:q-electron}
\end{equation}
Each new proton or neutron adds roughly $10^{-19}$ in Planck-scaled gravitational charge. Each electron contributes much less, but it is not zero.

\begin{table}[htbp]
\centering
\caption{Planck-scaled gravitational charge accounting for representative particles.}
\label{tab:particle-charge}
\begin{tabular}{lcc}
\toprule
Particle & $q_g=m/m_P$ & Self-coupling $q_g^2$ \\
\midrule
Electron & $4.18546\times 10^{-23}$ & $1.75181\times 10^{-45}$ \\
Proton & $7.68515\times 10^{-20}$ & $5.90615\times 10^{-39}$ \\
Neutron & $7.69574\times 10^{-20}$ & $5.92244\times 10^{-39}$ \\
Atomic mass unit $u$ & $7.62963\times 10^{-20}$ & $5.82113\times 10^{-39}$ \\
\bottomrule
\end{tabular}
\end{table}

\section{Periodic Table Confirmation Layer}

For an atom with atomic number $Z$, neutron number $N$, and mass number
\begin{equation}
 A=Z+N,
\label{eq:mass-number}
\end{equation}
the particle-count gravitational charge is
\begin{equation}
\boxed{
q_g(Z,N)=\frac{Zm_p+Nm_n+Zm_e-E_{\mathrm{bind}}/c^2}{m_P}.
}
\label{eq:isotope-qg}
\end{equation}
This is the isotope-level version. For an ordinary periodic-table element using average atomic weight $A_r(Z)$, the average atomic gravitational charge is
\begin{equation}
\boxed{\langle q_g\rangle_Z=A_r(Z)\frac{m_u}{m_P}.}
\label{eq:average-qg}
\end{equation}
This confirms the mass-counting part of the framework:
\begin{equation}
\boxed{\text{atomic gravitational charge rises with accumulated particle content}.}
\label{eq:charge-rises}
\end{equation}

\begin{table}[htbp]
\centering
\caption{Representative average atomic gravitational charges.}
\label{tab:atomic-charge}
\begin{tabular}{lccc}
\toprule
Element & $Z$ & Atomic weight or mass number & Average atomic $q_g$ \\
\midrule
H & 1 & 1.008 & $7.69067\times 10^{-20}$ \\
He & 2 & 4.002602 & $3.05384\times 10^{-19}$ \\
C & 6 & 12.011 & $9.16395\times 10^{-19}$ \\
O & 8 & 15.999 & $1.22066\times 10^{-18}$ \\
Fe & 26 & 55.845 & $4.26077\times 10^{-18}$ \\
Au & 79 & 196.96657 & $1.50278\times 10^{-17}$ \\
Pb & 82 & 207.2 & $1.58086\times 10^{-17}$ \\
U & 92 & 238.02891 & $1.81607\times 10^{-17}$ \\
Og & 118 & $[294]$ & $2.24311\times 10^{-17}$ \\
\bottomrule
\end{tabular}
\end{table}

The periodic-table layer does not, by itself, prove the Bessel-fold entropy engine. It confirms the accounting layer: more protons, neutrons, and electrons create more Planck-scaled gravitational source. The deeper claim requires a residual term that cannot be explained by ordinary rest mass and nuclear binding energy alone.

\section{Law of Large Numbers and Smooth Geometry}

For one particle,
\begin{equation}
 q_g=\frac{m}{m_P}.
\label{eq:q-single}
\end{equation}
For many particles,
\begin{equation}
\boxed{Q_g=\sum_{i=1}^{N}q_{g,i}.}
\label{eq:q-total}
\end{equation}
The average charge per particle is
\begin{equation}
 \bar q_g=\frac{1}{N}\sum_{i=1}^{N}q_{g,i}.
\label{eq:q-average}
\end{equation}
By the law of large numbers,
\begin{equation}
\boxed{\bar q_g\rightarrow\mathbb{E}[q_g]\quad\text{as}\quad N\rightarrow\infty.}
\label{eq:lln}
\end{equation}
The fluctuation shrinks as
\begin{equation}
\boxed{\sigma_{\bar q}\sim\frac{\sigma_q}{\sqrt{N}}.}
\label{eq:sigma-shrink}
\end{equation}
This is where gravity as accumulated force becomes mathematically strong:
\begin{equation}
\boxed{\text{single-particle gravitational charge is tiny and noisy},}
\label{eq:single-noisy}
\end{equation}
while
\begin{equation}
\boxed{\text{bulk accumulated gravitational charge becomes smooth and geometric}.}
\label{eq:bulk-smooth}
\end{equation}
The bridge is therefore
\begin{equation}
\boxed{
\text{particle counting}
\rightarrow
\text{large-number smoothing}
\rightarrow
\text{effective stress-energy}
\rightarrow
\text{curvature}.
}
\label{eq:counting-to-curvature}
\end{equation}

\section{Residual Prediction Beyond Ordinary Mass-Energy}

The final compact theory statement is the chain
\begin{align}
\boxed{x^2y''+xy'+(x^2-\nu^2)y=0,}\label{eq:compact-bessel}\\
\boxed{\nu=\ell+\frac{d-2}{2},\qquad d=4\Rightarrow\nu=\ell+1,}\label{eq:compact-order}\\
\boxed{\Delta E_{1/2}=\hbar v_s\frac{j_{\nu,q}-j_{\nu-1/2,q}}{R},}\label{eq:compact-energy}\\
\boxed{\Sigmahalf=\ell_P^2(\nabla_\mu\nabla_\nu-g_{\mu\nu}\Box)\left(\frac{\Delta S_{1/2}}{k_B}\right),}\label{eq:compact-sigma}\\
\boxed{\mathcal{G}_{\mu\nu}^{\mathrm{acc}}=\frac{8\pi G}{c^4}\int_{-\infty}^{t}\mathcal{K}(t,\tau)\left[\Tmat+\Tsound+\frac{c^4}{8\pi G\ell_P^2}\Sigmahalf\right]\dd\tau,}\label{eq:compact-gacc}\\
\boxed{q_g(X)=\frac{m_X}{m_P},}\label{eq:compact-qg}\\
\boxed{q_g(Z,N)=\frac{Zm_p+Nm_n+Zm_e-E_{\mathrm{bind}}/c^2}{m_P}.}\label{eq:compact-isotope}
\end{align}

Thus the refined claim is
\begin{equation}
\boxed{
\begin{minipage}{0.9\linewidth}
Mass-energy becomes geometry because resonant field energy, particle count, and entropy-producing Bessel-fold defects accumulate through a causal temporal kernel until their large-number average becomes smooth curvature.
\end{minipage}}
\label{eq:final-claim}
\end{equation}

The claim is therefore not split into a ``safe'' accounting piece and a disposable speculative piece. It is a single ladder of constraints. Planck-scaled gravitational charge gives the mass-accounting baseline; hyperspherical enclosure modes and Bessel roots define the resonant spectrum; ripple ratios and Pythagorean holonomy provide projection tests; and the entropy-fold term becomes meaningful precisely because it makes a residual prediction beyond ordinary mass-energy.

The clean residual target is
\begin{equation}
\boxed{
\Delta\mathcal{G}_{\mu\nu}^{\mathrm{res}}
=
\mathcal{G}_{\mu\nu}^{\mathrm{acc}}-
\mathcal{G}_{\mu\nu}^{\mathrm{matter}}
=
\frac{8\pi G}{c^4}\int_{-\infty}^{t}\mathcal{K}\Tsound\dd\tau
+
\frac{1}{\ell_P^2}\int_{-\infty}^{t}\mathcal{K}\Sigmahalf\dd\tau.
}
\label{eq:residual-target}
\end{equation}
If the ordinary rest-mass, binding-energy, and standard stress-energy contribution is subtracted and no residual remains, the Bessel-fold engine fails. If a reproducible residual remains with the predicted half-Bessel fingerprint, the model has its experimental opening.

\section{Conclusion}

The paper has assembled the proposed mechanism as a unified proof spine. Hyperspherical enclosure acoustics supplies the mode geometry. Bessel boundary roots supply the allowed spectrum. The half-order disharmonic fold supplies a concrete mode-drop fingerprint, approximately $0.180106$. Cross-scale ripple ratios supply the spatial projection test. Pythagorean comma holonomy supplies the frequency projection test. Planck-scaled mass accounting and large-number smoothing explain how tiny particle charges become smooth bulk curvature.

The key discipline is not to quarantine the new mechanism away from the established accounting, but to make each rung of the ladder answer to a different kind of check. The measurable question is whether the accumulated curvature/source contains a residual beyond ordinary mass-energy, and whether that residual carries the Bessel half-drop, projection-ratio, or holonomy signature predicted by the spine.

\appendix
\section{Formal Verification Scripts}
\label{app:formal-verification}

This appendix records the compact proof layer used to check the mathematical spine of the paper. The Wolfram Language checks verify the analytic and numerical reductions. The Isabelle layer records the logical dependency chain inside a bounded axiomatic framework, so that the computational certificate and the formal proof skeleton point to the same definitions used in the main text.

\subsection{Wolfram Language Certificate}

The following Wolfram Language checks evaluate the half-Bessel fingerprint, the causal kernel normalization, and the residual decomposition used in Sections 3 and 5.

\begin{lstlisting}
ClearAll["Global`*"];

nuD = ell + (d - 2)/2;
nuD4Residual = FullSimplify[(nuD /. d -> 4) - (ell + 1)];
nuD4Pass = TrueQ[nuD4Residual == 0];

j11 = BesselJZero[1, 1];
halfRoot = Pi;
dropFraction = N[(j11 - halfRoot)/j11, 16];
dropPass =
  TrueQ[0 < dropFraction < 1] &&
  Abs[dropFraction - 0.18010602117791018] < 10^-12;

deltaE = hbar vs (jnu - jhalf)/R;
deltaEZeroResidual = FullSimplify[deltaE /. jnu -> jhalf];
deltaEPass = TrueQ[deltaEZeroResidual == 0];
\end{lstlisting}

The causal kernel is normalized over past time by changing variables to $s=t-\tau$.

\begin{lstlisting}
kernelIntegral =
  Assuming[tdec > 0,
    Integrate[(1/tdec) Exp[-s/tdec], {s, 0, Infinity}]
  ];
kernelPass = TrueQ[kernelIntegral == 1];
\end{lstlisting}

The effective source and curvature residual decompose exactly into ordinary matter, resonance stress, and the Planck-scaled entropy-fold source.

\begin{lstlisting}
entropyStress = c^4/(8 Pi G ellP^2) Sigma;
residualSource = FullSimplify[(Tm + Ts + entropyStress) - Tm];
sourceDecompPass = TrueQ[residualSource == Ts + entropyStress];

curvMatter = (8 Pi G/c^4) IntK Tm;
curvTotal = (8 Pi G/c^4) IntK (Tm + Ts + entropyStress);
curvResidual = FullSimplify[curvTotal - curvMatter];
expectedCurvResidual =
  FullSimplify[
    (8 Pi G/c^4) IntK Ts + (1/ellP^2) IntK Sigma
  ];
curvPass =
  TrueQ[FullSimplify[curvResidual - expectedCurvResidual] == 0];
\end{lstlisting}

The Planck-scaled gravitational charge and isotope mass accounting checks are:

\begin{lstlisting}
qg = m/Sqrt[hbar c/G];
planckProductResidual =
  FullSimplify[
    (mX Sqrt[G/(hbar c)]) (mY Sqrt[G/(hbar c)])
      - G mX mY/(hbar c),
    Assumptions -> hbar > 0 && c > 0 && G > 0
  ];
planckPass = TrueQ[planckProductResidual == 0];

isotopeMass = Z mp + Nn mn + Z me - Ebind/c^2;
isotopeQg = isotopeMass/mP;
isotopeResidual =
  FullSimplify[mP isotopeQg - isotopeMass,
    Assumptions -> mP != 0 && c != 0
  ];
isotopePass = TrueQ[isotopeResidual == 0];
\end{lstlisting}

\subsection{Isabelle Concrete Mathematics Certificate}

The Isabelle layer has been strengthened beyond an abstract proof-graph certificate. The companion theory \texttt{TZPID\_HypersphericalBesselResidualBridge\_Math\_Checks.thy} uses typed real-valued definitions and structured Isar proofs for the algebraic pieces that Isabelle can check directly. The special-function numeric values, such as the first zero of $J_1$, remain in the Wolfram certificate; Isabelle proves the real-algebra identities that connect those values to the paper's proof spine.

\begin{lstlisting}
theory TZPID_HypersphericalBesselResidualBridge_Math_Checks
  imports Complex_Main
begin

definition hyperspherical_order :: "real => real => real" where
  "hyperspherical_order d ell = ell + (d - 2) / 2"

theorem hyperspherical_order_reduces_in_four_spatial_dimensions:
  "hyperspherical_order 4 ell = ell + 1"
proof -
  have "hyperspherical_order 4 ell = ell + (4 - 2) / 2"
    unfolding hyperspherical_order_def
    by (rule refl)
  also have "... = ell + 1"
    by linarith
  finally show ?thesis .
qed
\end{lstlisting}

The Pythagorean comma and its reciprocal bulk flip are proved as exact rational identities.

\begin{lstlisting}
definition pythagorean_comma :: real where
  "pythagorean_comma = (3 / 2) ^ 12 / 2 ^ 7"

definition pythagorean_bulk_reciprocal :: real where
  "pythagorean_bulk_reciprocal = 1 / pythagorean_comma"

theorem pythagorean_comma_exact:
  "pythagorean_comma = 531441 / 524288"
proof -
  have "pythagorean_comma = (3 / 2) ^ 12 / 2 ^ 7"
    unfolding pythagorean_comma_def
    by (rule refl)
  also have "... = 531441 / 524288"
    by norm_num
  finally show ?thesis .
qed

theorem pythagorean_comma_bulk_reciprocal_closes:
  "pythagorean_comma * pythagorean_bulk_reciprocal = 1"
proof -
  have "pythagorean_comma ~= 0"
  proof -
    have "pythagorean_comma = 531441 / 524288"
      using pythagorean_comma_exact .
    thus ?thesis
      by norm_num
  qed
  hence "pythagorean_comma * (1 / pythagorean_comma) = 1"
    by (rule nonzero_mult_div_cancel_left)
  thus ?thesis
    unfolding pythagorean_bulk_reciprocal_def .
qed
\end{lstlisting}

The ripple-index claim is also formalized: if wavelength and height are rescaled by the same nonzero projection factor, the dimensionless ripple index is unchanged.

\begin{lstlisting}
definition ripple_index :: "real => real => real" where
  "ripple_index wavelength height = wavelength / height"

theorem ripple_index_scale_invariant:
  assumes "scale ~= 0"
  shows "ripple_index (scale * wavelength) (scale * height)
       = ripple_index wavelength height"
proof -
  have "ripple_index (scale * wavelength) (scale * height)
      = (scale * wavelength) / (scale * height)"
    unfolding ripple_index_def
    by (rule refl)
  also have "... = wavelength / height"
    using assms
    by (field)
  finally show ?thesis
    unfolding ripple_index_def .
qed
\end{lstlisting}

Finally, Isabelle proves the Planck-charge product, isotope charge recovery, and residual-source decomposition as typed algebra.

\begin{lstlisting}
definition gravitational_charge :: "real => real => real" where
  "gravitational_charge mass mP = mass / mP"

theorem isotope_charge_recovers_isotope_mass:
  assumes "mP ~= 0"
  shows "mP * isotope_gravitational_charge Z N mp mn me Ebind c mP
       = isotope_mass Z N mp mn me Ebind c"
proof -
  have "mP * isotope_gravitational_charge Z N mp mn me Ebind c mP
      = mP * (isotope_mass Z N mp mn me Ebind c / mP)"
    unfolding isotope_gravitational_charge_def
    by (rule refl)
  also have "... = isotope_mass Z N mp mn me Ebind c"
    using assms
    by (field)
  finally show ?thesis .
qed

definition residual_curvature :: "real => real => real => real" where
  "residual_curvature total matter = total - matter"

definition accumulated_total_source :: "real => real => real => real" where
  "accumulated_total_source matter sound entropy =
     matter + sound + entropy"

theorem residual_curvature_decomposes:
  "residual_curvature
     (accumulated_total_source matter sound entropy) matter
   = sound + entropy"
proof -
  have "residual_curvature
          (accumulated_total_source matter sound entropy) matter
      = (matter + sound + entropy) - matter"
    unfolding residual_curvature_def accumulated_total_source_def
    by (rule refl)
  also have "... = sound + entropy"
    by algebra
  finally show ?thesis .
qed
\end{lstlisting}

\subsection{Expanded Phase 2 HOL Coverage}

The Phase 2 Isabelle layer has also been widened beyond the twelve Bessel-residual bridge obligations. The theory \texttt{TZPID\_Phase2\_Expanded\_Theorem\_Coverage.thy} records the paper-facing theorem neighborhood as eight connected families with seventy-eight target-ID mentions: the Hyperspherical Bessel residual bridge, nested hyperspherical enclosure, cosmic acoustics, cross-scale ripple projection, Pythagorean/Hopf reciprocal holonomy, critical scale-invariance, gyromagnetic movement, and Kuramoto-style phase-locking/orbital synchronization. This does not claim that every theorem row in the full registry has already been semantically translated into HOL; it marks the expanded formal scope used by this paper.

\begin{lstlisting}
theorem expanded_phase2_spine_connections:
  "hyperspherical_bessel_residual_bridge_chain
   & cosmic_acoustics_chain
   & ripple_projection_chain
   & acoustic_holonomy_chain
   & critical_scale_invariance_chain
   & nested_hypersphere_unifying_chain
   & gyromagnetic_movement_chain
   & phase_locking_resonance_chain"
proof (intro conjI)
  show "hyperspherical_bessel_residual_bridge_chain" using tap_chain .
  show "cosmic_acoustics_chain" using cosmic_chain .
  show "ripple_projection_chain" using ripple_chain .
  show "acoustic_holonomy_chain" using acoustic_holonomy_chain .
  show "critical_scale_invariance_chain" using critical_chain .
  show "nested_hypersphere_unifying_chain" using unifying_chain .
  show "gyromagnetic_movement_chain" using movement_chain .
  show "phase_locking_resonance_chain" using ratio_chain .
qed
\end{lstlisting}

The Kuramoto-style orbital synchronization mechanism is represented as typed HOL algebra in \texttt{TZPID\_HypersphericalBesselResidualBridge\_Phase2\_Model.thy}. Isabelle proves both zero-drift phase lock for equal intrinsic frequencies and the rational $3{:}2$ orbital lock used by the resonance-selection spine.

\begin{lstlisting}
definition kuramoto_pair_drift :: "real => real => real => real => real" where
  "kuramoto_pair_drift omega_i omega_j coupling phase_gap =
     omega_i - omega_j - coupling * sin phase_gap"

definition rational_orbital_lock :: "real => real => real => real => bool" where
  "rational_orbital_lock p q spin_frequency orbit_frequency =
     (q * spin_frequency = p * orbit_frequency)"

theorem kuramoto_equal_frequency_zero_phase_locks:
  "kuramoto_pair_drift omega omega coupling 0 = 0"

theorem three_two_orbital_lock:
  "rational_orbital_lock 3 2 ((3 / 2) * orbit_frequency) orbit_frequency"
\end{lstlisting}

The full Isabelle session was rebuilt after adding these theories. The normalization integral for the causal exponential kernel remains certified in Wolfram until the Isabelle side imports and develops the required HOL-Analysis integration proof.

\section*{Acknowledgements}
\addcontentsline{toc}{section}{Acknowledgements}

Thank you to Codex, Claude, and Gemini for their assistance in developing, testing, and refining the conceptual structure of this paper.

\begin{thebibliography}{9}
\bibitem{dlmf-bessel}
NIST Digital Library of Mathematical Functions, Section 10.2, ``Definitions,'' Bessel and Hankel Functions. Available at \url{https://dlmf.nist.gov/10.2}.

\bibitem{codata2022}
E. Tiesinga, P. J. Mohr, D. B. Newell, and B. N. Taylor, ``CODATA Recommended Values of the Fundamental Physical Constants: 2022,'' arXiv:2409.03787. Available at \url{https://arxiv.org/abs/2409.03787}.

\bibitem{iupac-periodic}
International Union of Pure and Applied Chemistry, ``Periodic Table of Elements.'' Available at \url{https://iupac.org/what-we-do/periodic-table-of-elements/}.
\end{thebibliography}

\end{document}

